
%(BEGIN_QUESTION)
% Copyright 2006, Tony R. Kuphaldt, released under the Creative Commons Attribution License (v 1.0)
% This means you may do almost anything with this work of mine, so long as you give me proper credit

På kalde dager oppleves temperaturen lavere når det blåser. Selv om et termometer ikke viser lavere temperatur når vinden øker, føles det kaldere. Denne effekten brukes i værmeldinger som ``følt temperatur''.

Den samme effekten kan brukes til å måle massestrøm. Strømningsmålere basert på dette prinsippet kalles termiske strømningsmålere. Forklar hvordan en termisk strømningsmåler er konstruert, og angi hvilke fysiske egenskaper ved fluidstrømmen som påvirker kalibreringen.

\underbar{file i00534}
%(END_QUESTION)





%(BEGIN_ANSWER)

Termiske massestrømningsmålere bruker et elektrisk varmeelement og minst to temperatursensorer for å registrere temperaturforskjell knyttet til konveksjon.

Kalibreringen av en termisk måler avhenger av fluidets varmeledningsevne, samt dets spesifikke varmekapasitet (hvor mye varmeenergi som tas opp per masseenhet per temperaturøkning), på samme måte som kalibreringen av en orifiseplate avhenger av fluidets tetthet. Hvis én eller begge disse faktorene varierer i prosessstrømmen, vil en termisk massestrømningsmåler gi ustabile indikasjoner, akkurat som en orifiseplate gir ustabile avlesninger når prosessfluidets tetthet endrer seg tilfeldig.

\vskip 10pt

Et eksempel på et stoff med svært høy spesifikk varmekapasitet er \textit{hydrogen}gass. Hvis en termisk massestrømningsmåler er kalibrert for luft, og deretter eksponeres for en strøm av hydrogengass, vil den feilaktig indikere for høy strømning av hydrogen på grunn av gassens svært høye spesifikke varmekapasitet.

%(END_ANSWER)





%(BEGIN_NOTES)

%INDEX% Måling, strømning: termisk (masse)

%(END_NOTES)

